
\section{Introduction}
The increasing deployment of Ethernet-based infrastructures has weakened the assumption that the data-link layer is inherently trustworthy. In these settings, link-layer traffic is exposed to adversarial conditions, and cryptographic protection is needed to preserve confidentiality, integrity, and replay resistance. The IEEE 802.1AE MACsec standard~\cite{ieee8021ae} addresses this need by defining hop-by-hop protection for Ethernet frames, while the MACsec Key Agreement protocol (MKA)~\cite{ieee8021x} manages the establishment and maintenance of the cryptographic state used by MACsec.

MKA operates within a Connectivity Association (CA), where an elected Key Server derives fresh Secure Association Keys (SAKs) and distributes them to participating peers. The protocol relies on a pre-shared Connectivity Association Key (CAK), from which derived keys such as the Key Encryption Key (KEK) and the Integrity Check Value Key (ICK) are computed. Message Numbers (MNs) and Key Numbers (KNs) are used to track replay protection and key-generation epochs, including periodic rekeying. Despite the practical importance of MKA, formal analyses of its key-agreement logic remain limited. Existing work has largely focused on MACsec data-plane mechanisms or performance aspects rather than on machine-checked verification of the protocol's core security guarantees.

This paper addresses that gap by presenting a formal analysis of a simplified two-party MKA exchange using the Tamarin prover. Our model captures the initial session establishment and a subsequent rekey round, and we verify essential security properties under a standard Dolev--Yao adversary model. The analysis confirms the expected guarantees when the CAK remains uncompromised, while also revealing a structural weakness in the protocol's acceptance logic.

\paragraph{Contributions.} This work makes three contributions. First, it constructs a Tamarin model of a two-party MKA exchange that captures the initial session and a rekey round, including the monotonic MN/KN counters. Second, it formalizes and verifies core security properties such as executability, agreement, authentication, secrecy, ordering, and freshness. Third, it identifies and demonstrates a structural weakness in the protocol's acceptance logic: a malicious or compromised Key Server that possesses the CAK can inject an arbitrary SAK that the Server accepts.

\paragraph{Organization.} Section~\ref{sec:background} explains MKA and its lifecycle. Section~\ref{sec:analysis} introduces the formal analysis and the Tamarin prover. Section~\ref{sec:model} presents our formal model and the protocol flow in Figure~\ref{fig:mka_flow}. Section~\ref{sec:results} reports the verification results, and Section~\ref{sec:conclusion} concludes the paper.
